
\subsubsection{XXX 模型的检验}

% 本节对所建立的【模型名称】进行检验。
% 检验内容:
%   1. 5 折交叉验证(评估泛化能力)
%   2. 误差分析(RMSE、MAE、MAPE 等)
%   3. 灵敏度分析(可选,关键参数扰动)
%   4. 稳健性检验(可选,替换算法或参数区间)
% 整体检验在第 6 章「六、模型检验」展开(误差/灵敏度/稳健性),此处为问题内的局部检验。

为验证本文所建立的【模型名称】的可行性和准确性,本节进行如下检验。

\textbf{(1) 5 折交叉验证。}采用 5 折交叉验证评估模型的泛化能力,【结果说明】。结果显示【指标】为【值 $\pm$ 标准差】,模型【无明显过拟合 / 存在过拟合风险】。

\textbf{(2) 误差分析。}采用 RMSE、MAE、MAPE 等指标量化模型精度,【结果说明】。相较【对比方法】,本文模型误差【大幅降低/相当】,拟合与计算精度【优异/可接受】。

\begin{longtable}{>{\centering\arraybackslash}p{0.22\textwidth}>{\centering\arraybackslash}p{0.22\textwidth}>{\centering\arraybackslash}p{0.22\textwidth}>{\centering\arraybackslash}p{0.22\textwidth}}
  \caption{问题 1 模型误差分析}\\
  \toprule
  指标 & 训练集 & 测试集 & 说明 \\
  \midrule
  \endfirsthead
  \caption{问题 1 模型误差分析(续)}\\
  \toprule
  指标 & 训练集 & 测试集 & 说明 \\
  \midrule
  \endhead
  \bottomrule
  \endfoot
  $R^2$ & 【值】 & 【值】 & 决定系数 \\
  RMSE & 【值】 & 【值】 & 均方根误差 \\
  MAE  & 【值】 & 【值】 & 平均绝对误差 \\
  MAPE(\%) & 【值】 & 【值】 & 平均绝对百分比误差 \\
\end{longtable}

\textbf{(3) 灵敏度分析。}对核心参数 $\alpha$、$\beta$ 等进行 $\pm 10\%$、$\pm 20\%$ 扰动测试,【结果说明】。参数小幅波动时,模型求解结果偏差控制在 5\% 以内,模型【抗干扰能力强 / 对 XX 参数敏感】。

\textbf{(4) 稳健性检验。}替换【对比方法】重新求解,核心结论趋势一致,误差可控,证明模型结论无偶然性。

综合上述检验,【模型名称】在【评估指标】上表现稳定,验证了模型的有效性。后续整体检验详见第 6 章。
